\scanpage{70}
% archive foliation: 69

Supposons $A$ \deleted{\unsure{Ded.}} Dedekind \unsure{et} local
\inserted{(cas Brauer)} de sorte que
\[
\operatorname{Is}(C)=1,\qquad\operatorname{Brd}(C)=\mathbb Z/e\mathbb Z
\]
\inserted{$\operatorname{Pic}(A)=0$}
\margin{$e$: ramification}
et on trouve
\[
\operatorname{Aut\,ext}(O)=\operatorname{Brd}C=\mathbb Z/e\mathbb Z
\]
Donc il y a obstruction locale à ce qu'un autom. de $O$ soit néc. intérieur, il y a $e$ classes d'autom de $O$.

Supposons \inserted{dans cas global} qu'un autom de $O$ soit \unsure{localmt} intérieur (pour la top. fppf par exemple)
\inserted{i.e. son image dans $\operatorname{Ram}(C/A)$ nulle},
\deleted{Alors il est défini}
avec une section de $\uline{O}^*/G_{m\,A}$, et \deleted{\unsure{elle a}}
\inserted{l'obstruction}
pour le relever est un élément de $\operatorname{Pic}(A)$, nulle dans
\deleted{$\operatorname{Pic}(O)$} $H^1(A,O^*)$
$\ldots$ donnant un élément trivial de $\operatorname{Is}(C)$ $\ldots$
\deleted{\unsure{Donc}}
ce qui se voyait aussi par les \uline{suites exactes} précédentes.

Cas des corps gauches de Deuring \unsure{splitté} en $p$ et $\infty$: la suite exacte devient
\diagram{\begin{tikzcd}[column sep=small,row sep=small]
0\arrow[r,dash]&\mathbb Z/2\mathbb Z\arrow[r]&O^*\arrow[r]\arrow[d,equal]&
\operatorname{Aut}(O)\arrow[r]&\mathbb Z/2\mathbb Z\arrow[r]&
\pi_0(C,X)\arrow[r]&\pi_0(\mathcal O,O)\arrow[r,dash]&1\\
&\begin{gathered}\text{engendré}\\\text{par }-1\end{gathered}&
\begin{gathered}\mathbb Z/2\mathbb Z\\
\text{sauf si}\\j=0\text{ ou }j=12^3\\
(\text{où on trouve}\\\mathbb Z/4\mathbb Z\text{ resp }\mathbb Z/6\mathbb Z\\
\text{si }p\ne2,3)\end{gathered}
&&\begin{gathered}\text{engendré}\\\text{par Frobenius}\end{gathered}&
\begin{gathered}
\text{classes d'}\\\text{is. de classes}\\
\text{de courbes}\\\text{elliptiques}\\
\text{de Hasse nul}\\
\text{l'opération de ``Frobenius''}\\
\text{dessus est }E\longmapsto E^{(p)}
\end{gathered}&
\begin{gathered}\text{classes de}\\\text{conjugaison}\\
\text{\unsure{possibles}}\\\text{d'ordres}\\\text{maximaux}\end{gathered}
\end{tikzcd}}
\margin{\unsure{on précisera} pour quels $p$ ces $j$ sont d'inv. de Hasse nul $\ldots$}
\margin{\unsure{épim} ssi la courbe elliptique $X$ d'anneau $O$ se descend à $\mathbb F_p$, i.e. son $j\in\mathbb F_p$}

Question: $O^*$ est-il tjrs égal
